\twocolumn[
\begin{@twocolumnfalse}
\begin{center}
{\large\bfseries Aerosol-generating steps in routine clinical laboratory testing:\\[2pt]
a procedure inventory and evidence brief\par}
\vspace{10pt}
{\normalsize Oyekunle Adedamola Oyeniyi\,\href{https://orcid.org/0000-0002-1978-6133}{\textcolor{linknavy}{$\circ$}}\par}
\vspace{3pt}
{\small\itshape Independent Researcher, Pullman, Washington 99163, USA\par}
\vspace{2pt}
{\footnotesize ORCID: \href{https://orcid.org/0000-0002-1978-6133}{0000-0002-1978-6133}
\quad$\cdot$\quad \href{mailto:adedamolaoyekunleniyi@gmail.com}{adedamolaoyekunleniyi@gmail.com}\par}
\vspace{4pt}
{\footnotesize Technical report \quad$\cdot$\quad Derived from Clinical Laboratory Aerosol Biorisk Reference v1.1\par}
\vspace{14pt}
\end{center}

\begin{adjustwidth}{0.055\textwidth}{0.055\textwidth}
\noindent\small
Routine diagnostic testing requires a laboratory worker to open, spin, pipette,
smear and subculture human specimens, and each of those actions is capable of
putting respirable material into the air. National biosafety guidance obliges
every laboratory to assess the risk of its own procedures, but no public
inventory sets out what the bench steps of routine testing actually are, what
is known about aerosol release from each, and how strong that knowledge is.
This brief supplies one. Seven bench steps common to moderate- and
high-complexity testing are defined by their boundaries, characterised for
release, and graded for evidence on a five-level scale in which a grade is
reported beside a score and never combined into it. The brief then assembles,
on a single axis, every airborne concentration measurement located for clinical
bench work. That assembly is its principal contribution, and it is short:
four conditions in which nothing was recovered above a detection limit of
4~CFU\,m$^{-3}$, two discrete values of 16 and 325~CFU\,m$^{-3}$ for
acid-fast smear preparation, and three simulated accident intensities spanning
roughly 10 to $10^{4}$~CFU\,m$^{-3}$. Three of seven steps carry direct
measurement, two rest on incident series with no aerosol measurement of either
step in a clinical laboratory, one rests on an inference from a research
setting, and one, the operation of automated chemistry and haematology
analysers, has no located evidence of any kind despite being the
highest-throughput step in the building. For tuberculosis work the measured
record is unusually favourable and unusually narrow: cartridge-based molecular
testing produced no culturable bioaerosol at any stage, against 16 and
325~CFU\,m$^{-3}$ for the smear microscopy it displaces, while every other step
on the pulmonary tuberculosis pathway scores Elevated or higher against a
low-dose inhalational profile on evidence that is either absent or
incident-based. Fourteen survey items are derived from the gaps, for use by
programmes surveying tuberculosis laboratory practice. The inventory, the
extracted evidence, the grading and the analysis code are released openly.

\vspace{7pt}
\noindent\small\textbf{Keywords:} aerosol-generating procedures; clinical
laboratory; biosafety; tuberculosis; GeneXpert; evidence grading; laboratory
survey; biorisk management
\vspace{16pt}
\end{adjustwidth}
\end{@twocolumnfalse}
]
